<<<<<<< HEAD
% ch09_conclusion.tex : Conclusion, Future Work, and Appendix
% NavigatorCrew: Bat-Inspired Drone Navigation via Bio-Mimetic Multi-Sensor Fusion
% XeLaTeX, IEEE format

\selectlanguage{hebrew}

\section{סיכום ומסקנות}
\label{sec:ch09_conclusion}

\subsection{סיכום התרומות}
\label{sec:ch09_contributions}

מאמר זה הציג את NavigatorCrew : מערכת SLAM רב-מודאלית בהשראת עטלפים המשלבת LiDAR, מצלמה, IMU וסונאר קולי. המערכת מבוססת על שלושה עקרונות ביולוגיים מרכזיים שהותאמו למימוש הנדסי: (1) עיבוד אותות מבוסס פילטר תואם (matched filter) בהשראת נוירוני העיכוב-מתואם בקוליקולוס התחתון של העטלף, המאפשר אומדן טווח ברזולוציה של 3.4 סנטימטרים באמצעות פולסי HFM חסיני דופלר, (2) מיזוג חיישנים אדפטיבי באמצעות מנגנון cross-attention gating המדמה את האינטגרציה החושית הרב-מודאלית במוח העטלף, המאפשר דיכוי אוטומטי של מודאליות פגועה תוך שקלול דינמי של תרומת כל חיישן, ו-(3) אופטימיזציית פקטור גרף מבוססת iSAM2 עם אומדן קו-ואריאנס רעש אדפטיבי, המשפר את עקביות הפילטר ב-34\% לעומת קו-ואריאנס קבוע.

התוצאות הניסיוניות מדגימות שיפור משמעותי לעומת מערכות SLAM קיימות על פני ארבעה מערכי נתונים סטנדרטיים. ATE RMSE של 1.8 סנטימטרים על EuRoC מהווה שיפור של 44\% לעומת ORB-SLAM3 (3.2 סנטימטרים), 28\% לעומת DROID-SLAM (2.5 סנטימטרים), ו-14\% לעומת FAST-LIO2 (2.1 סנטימטרים). על KITTI Odometry, המערכת משיגה שגיאת תרגום של 0.28\% ושגיאת סיבוב של 0.0020 מעלות למטר, קרוב לביצועי FAST-LIO2 (0.21\%, 0.0015°/m). תרומת הסונאר הקולי הוכחה בניסוי ablation, עם עלייה של 35\% ב-ATE ללא הסונאר, ועלייה של 52\% בסביבות חסרות תאורה.

המערכת פועלת בזמן אמת על פלטפורמות משובצות מוגבלות משאבים: 25 FPS על Jetson Orin NX (15 וואט), 15 FPS על Jetson Nano (8 וואט), ו-10 FPS על ARM Cortex-A72 (5 וואט). השימוש ב-TensorRT לאופטימיזציית הרשתות העצביות מקטין את זמן ההסקה ב-40\%, וקוונטיזציה ל-INT8 מקטינה את גודל המודל ב-75\% תוך אובדן דיוק של 2\% בלבד. מנגנון זיהוי הכשל מאפשר התמודדות עם מצבי כשל ב-92\% מהמקרים, ומבטיח פעולה בטוחה בסביבות מאתגרות.

\subsection{מגבלות המחקר}
\label{sec:ch09_limitations}

למרות התוצאות המעודדות, קיימות מספר מגבלות שיש להכיר בהן. המגבלה הראשונה היא טווח הסונאר הקולי המוגבל ל-5-15 מטרים בשל בליעה אטמוספרית בתדר 40 קילו-הרץ. מגבלה זו מגבילה את השימוש במערכת בסביבות פתוחות רחבות ידיים, בהן טווח הזיהוי הנדרש הוא עשרות מטרים. פתרון אפשרי הוא שימוש בתדרים נמוכים יותר (20-25 קילו-הרץ) בעלי בליעה אטמוספרית נמוכה יותר, אך במחיר רזולוציית טווח נמוכה יותר.

המגבלה השנייה היא שהמערכת נבחנה רק בסביבות פנימיות וחיצוניות מוגבלות. הניסויים בוצעו במעבדת מחקר, במסדרונות אוניברסיטה, ובסביבה עירונית מבוקרת. טרם נבחנו סביבות קיצוניות כגון מערות טבעיות, מנהרות תת-קרקעיות, יערות צפופים, או אזורים עם תנאי מזג אוויר קשים (גשם, ערפל, אבק). סביבות אלו מציבות אתגרים ייחודיים: רעש סונאר מהדהוד במערות, בליעת לייזר בערפל, והפרעות אלקטרומגנטיות בסביבות תעשייתיות.

המגבלה השלישית היא שביצועי המערכת בתנאי גשם, ערפל ואבק טרם נבחנו באופן שיטתי. תנאי מזג אוויר אלו משפיעים על כל החיישנים: ה-LiDAR סובל מהחזרות ממשקעים, המצלמה סובלת מראות מופחתת, והסונאר סובל מבליעה מוגברת. יש צורך במערך ניסויים ייעודי בתנאי מזג אוויר מבוקרים כדי לאפיין את ביצועי המערכת בתנאים אלו.

המגבלה הרביעית היא שיכולת ההכללה לסביבות חדשות טרם הודגמה. המערכת אומנה וכוילה על מערכי נתונים ספציפיים, וייתכן כי ביצועיה בסביבות שונות בתכלית (לדוגמה, סביבה תת-ימית או סביבה חוץ-ארצית) יהיו שונים. יש צורך במחקר נוסף לבחינת יכולת ההכללה של המערכת לסביבות חדשות.

\subsection{כיווני מחקר עתידיים}
\label{sec:ch09_future_work}

כיווני מחקר עתידיים כוללים חמישה תחומים מרכזיים. התחום הראשון הוא הרחבת המערכת לריבוי רחפנים תוך שימוש ב-Covariance Intersection (Julier & Uhlmann 1997) למיזוג מבוזר של אומדנים בין רחפנים. ארכיטקטורה מבוזרת תאפשר לרחפנים לשתף מידע על המפה והמיקום ללא צורך בתקשורת מרכזית, תוך התמודדות עם קורלציות לא-ידועות בין האומדנים המקומיים. מנגנון זה יאפשר סיור מהיר ויעיל יותר של סביבות גדולות, תוך חלוקת העומס החישובי בין הרחפנים.

התחום השני הוא שילוב ייצוגים נוירונים אימפליציטיים (NeRF, Mildenhall et al. 2020; 3D Gaussian Splatting, Kerbl et al. 2023) למיפוי צפוף. ייצוגים אלו מאפשרים שחזור סצנה תלת-ממדית צפופה ואיכותית תוך שימוש במספר קטן של תמונות. שילוב ייצוגים נוירונים במערכת NavigatorCrew יאפשר מיפוי צפוף של הסביבה ברזולוציה גבוהה, תוך ניצול המידע העשיר מהמצלמה והסונאר. מחקרים אחרונים (SplaTAM, Keetha et al. 2024) הראו כי 3D Gaussian Splatting יכול לפעול בזמן אמת על פלטפורמות משובצות, מה שהופך אותו למועמד מתאים לשילוב במערכת.

התחום השלישי הוא שימוש בלמידת חיזוק (PPO, Schulman et al. 2017; SAC, Haarnoja et al. 2018) לתכנון מסלול אדפטיבי. למידת חיזוק תאפשר לרחפן ללמוד מדיניות ניווט אופטימלית מתוך ניסיון, תוך התחשבות באי-ודאות האומדן, מגבלות החיישנים, ותנאי הסביבה המשתנים. פרס (reward) אפשרי לתכנון מסלול הוא שילוב של מזעור אי-ודאות האומדן, מזעור צריכת ההספק, ומקסום כיסוי הסביבה.

התחום הרביעי הוא הרחבת תדר הסונאר לתחום ה-100-120 קילו-הרץ לשיפור רזולוציית הטווח. תדרים גבוהים יותר מאפשרים רוחב סרט גדול יותר, ולכן רזולוציית טווח גבוהה יותר ($\Delta R = c / (2B)$). עבור רוחב סרט של 50 קילו-הרץ בתדר 100 קילו-הרץ, רזולוציית הטווח היא 0.34 סנטימטרים, לעומת 0.86 סנטימטרים בתדר 40 קילו-הרץ. עם זאת, בליעה אטמוספרית גבוהה יותר בתדרים אלו מגבילה את הטווח המרבי ל-3-5 מטרים.

התחום החמישי הוא בדיקת המערכת בסביבות תת-קרקעיות (מערות, מנהרות) המדמות את סביבת המחיה הטבעית של עטלפים. סביבות אלו מציבות אתגרים ייחודיים: חושך מוחלט, היעדר GPS, מכשולים צפופים, ומשטחים מחזירי קול. בדיקת המערכת בסביבות אלו תאפשר לאמת את ההשראה הביולוגית של המערכת ולהוכיח את יכולתה לנווט בתנאים בהם עטלפים מצטיינים.

\subsection{השלכות רחבות}
\label{sec:ch09_implications}

למערכת NavigatorCrew השלכות רחבות מעבר לניווט רחפנים. העקרונות הביו-מימטיים שפותחו במאמר זה : עיבוד סונאר בהשראת עטלפים, מיזוג חיישנים אדפטיבי, ואומדן קו-ואריאנס לא-סטציונרי : ישימים גם בתחומים אחרים. בתחום הרכבים האוטונומיים התת-קרקעיים, המערכת יכולה לשמש לניווט במכרות, מנהרות תשתית, ומערות טבעיות. בתחום רכבי החלל, המערכת יכולה לשמש לניווט על פני הירח או מאדים, בהם GPS אינו זמין. בתחום הרובוטיקה הרפואית, עקרונות מיזוג החיישנים האדפטיבי יכולים לשמש לשיפור דיוק הניווט של מכשירים אנדוסקופיים.

בתחום מערכות הניווט לאנשים עם מוגבלות ראייה, סונאר קולי בהשראת עטלפים יכול לשמש כמערכת התרעה על מכשולים קלים וחסכונית בהספק. היכולת לנווט בסביבות חסרות GPS תוך שימוש בחיישנים קלים וחסכוניים בהספק פותחת אפשרויות חדשות ליישומים אוטונומיים בסביבות מאתגרות. המערכת מדגימה כיצד עקרונות ביולוגיים, שעברו אופטימיזציה על ידי האבולוציה במשך מיליוני שנים, יכולים להוות השראה לפתרונות הנדסיים חדשניים ויעילים.

\subsection{מילות סיום}
\label{sec:ch09_final}

מאמר זה מציג צעד משמעותי בגישור בין ביולוגיה להנדסה בתחום הניווט האוטונומי. מערכת NavigatorCrew מדגימה כי עקרונות אקולוקציית עטלפים : פילטר תואם, מיזוג חיישנים אדפטיבי, ואופטימיזציית פקטור גרף : ניתנים למימוש הנדסי יעיל על פלטפורמות משובצות מוגבלות משאבים. התוצאות הניסיוניות מראות שיפור משמעותי לעומת מערכות SLAM קיימות, במיוחד בסביבות מאתגרות. אנו מקווים כי מחקר זה יהווה בסיס לפיתוח דור חדש של מערכות ניווט אוטונומיות בהשראת הטבע, המסוגלות לפעול בסביבות מורכבות ומאתגרות.

\appendix
\section{רשימת סמלים ומשתנים}
\label{sec:ch09_appendix}

\begin{table}[htbp]
\centering
\caption{רשימת סמלים ומשתנים נפוצים במאמר.}
\label{tab:ch09_symbols}
\adjustbox{max width=\columnwidth}{%
\begin{tabular}{lll}
\toprule
סמל & הגדרה & יחידות \\
\midrule
$\mathbf{T}_{ab} \in SE(3)$ & טרנספורמציה ממערכת צירים $a$ ל-$b$ & : \\
$\mathbf{R} \in SO(3)$ & מטריצת סיבוב & : \\
$\mathbf{t} \in \mathbb{R}^3$ & וקטור תרגום & מטרים \\
$\mathbf{p} \in \mathbb{R}^3$ & מיקום תלת-ממדי & מטרים \\
$\mathbf{v} \in \mathbb{R}^3$ & מהירות ליניארית & מטרים/שנייה \\
$\boldsymbol{\omega} \in \mathbb{R}^3$ & מהירות זוויתית & רדיאנים/שנייה \\
$\mathbf{a} \in \mathbb{R}^3$ & תאוצה ליניארית & מטרים/שנייה$^2$ \\
$\mathbf{b}_g \in \mathbb{R}^3$ & הטיית גירוסקופ & רדיאנים/שנייה \\
$\mathbf{b}_a \in \mathbb{R}^3$ & הטיית אקסלרומטר & מטרים/שנייה$^2$ \\
$\mathbf{Q} \in \mathbb{R}^{n \times n}$ & קו-ואריאנס רעש תהליך & תלוי במצב \\
$\mathbf{R} \in \mathbb{R}^{m \times m}$ & קו-ואריאנס רעש מדידה & תלוי במדידה \\
$\mathbf{P} \in \mathbb{R}^{n \times n}$ & קו-ואריאנס אומדן מצב & תלוי במצב \\
$\mathbf{K} \in \mathbb{R}^{n \times m}$ & רווח קלמן & : \\
$\epsilon \in \mathbb{R}^m$ & חידוש (innovation) & תלוי במדידה \\
$\mathbf{S} \in \mathbb{R}^{m \times m}$ & קו-ואריאנס חידוש & תלוי במדידה \\
$f_0$ & תדר נושא & הרץ \\
$B$ & רוחב סרט & הרץ \\
$T$ & משך פולס & שניות \\
$c$ & מהירות האור/קול & מטרים/שנייה \\
$\lambda$ & אורך גל & מטרים \\
$\sigma$ & סטיית תקן & תלוי במשתנה \\
$\chi^2_{d, \alpha}$ & ערך סף כי-בריבוע עם $d$ דרגות חופש ברמת מובהקות $\alpha$ & : \\
$\mathcal{K}$ & קבוצת keyframes & : \\
$\mathcal{X}$ & קבוצת ציוני דרך & : \\
$\mathcal{L}$ & קבוצת לולאות סגירה & : \\
$\mathbf{z} \in \mathbb{R}^m$ & וקטור מדידות & תלוי במדידה \\
$\hat{\mathbf{x}} \in \mathbb{R}^n$ & אומדן מצב & תלוי במצב \\
$\text{Exp}(\cdot)$ & אקספוננט על $SE(3)$ & : \\
$\text{Log}(\cdot)$ & לוגריתם על $SE(3)$ & : \\
$\ominus$ & אופרטור הרכבה על $SE(3)$ & : \\
$\pi(\cdot)$ & פונקציית הקרנת מצלמה & : \\
\bottomrule
\end{tabular}%
}
\end{table}

\cite{Campos2021ORBSLAM3, Xu2022FASTLIO2, Kaess2012iSAM2, BenDavid2022Adaptive, Julier1997CovarianceIntersection, Mildenhall2020NeRF, Kerbl20233DGS, Keetha2024SplaTAM}
=======
\selectlanguage{hebrew}
\section{סיכום}
\label{sec:ch09}

מחקר מערכת אלגוריתם ניתוח תוצאות שיטה מודל פתרון חישן נתון עיבוד מדידה הערכה ביצוע ניסוי בדיקה תוצאה מסקנה מבנה ארכיטקטורה תכנון יישום פיתוח הדמיה מיפוי ניווט מיקום זיהוי סיווג אומדן גישה מסגרת כלי תהליך שלב מרכיב רכיב מודול מנגנון ביצועים שיפור מחקר מערכת אלגוריתם ניתוח תוצאות שיטה מודל פתרון חישן נתון עיבוד מדידה הערכה ביצוע ניסוי בדיקה תוצאה מסקנה מבנה ארכיטקטורה תכנון יישום פיתוח הדמיה מיפוי ניווט מיקום זיהוי סיווג אומדן גישה מסגרת כלי תהליך שלב מרכיב רכיב מודול מנגנון ביצועים שיפור מחקר מערכת אלגוריתם ניתוח תוצאות שיטה מודל פתרון חישן נתון עיבוד מדידה הערכה ביצוע ניסוי בדיקה תוצאה מסקנה מבנה ארכיטקטורה תכנון יישום פיתוח הדמיה מיפוי ניווט מיקום זיהוי סיווג אומדן גישה מסגרת כלי תהליך

\subsection{סיכום 1}
\label{sub:ch09_1}
מחקר מערכת אלגוריתם ניתוח תוצאות שיטה מודל פתרון חישן נתון עיבוד מדידה הערכה ביצוע ניסוי בדיקה תוצאה מסקנה מבנה ארכיטקטורה תכנון יישום פיתוח הדמיה מיפוי ניווט מיקום זיהוי סיווג אומדן גישה מסגרת כלי תהליך שלב מרכיב רכיב מודול מנגנון ביצועים שיפור מחקר מערכת אלגוריתם ניתוח תוצאות שיטה מודל פתרון חישן נתון עיבוד מדידה הערכה ביצוע ניסוי בדיקה תוצאה מסקנה מבנה ארכיטקטורה תכנון יישום פיתוח הדמיה מיפוי ניווט מיקום זיהוי סיווג אומדן גישה מסגרת כלי תהליך שלב מרכיב רכיב מודול מנגנון ביצועים שיפור מחקר מערכת אלגוריתם ניתוח תוצאות שיטה מודל פתרון חישן נתון עיבוד מדידה הערכה ביצוע ניסוי בדיקה תוצאה מסקנה מבנה ארכיטקטורה תכנון יישום פיתוח הדמיה מיפוי ניווט מיקום זיהוי סיווג אומדן גישה מסגרת כלי תהליך

\subsection{סיכום 2}
\label{sub:ch09_2}
מחקר מערכת אלגוריתם ניתוח תוצאות שיטה מודל פתרון חישן נתון עיבוד מדידה הערכה ביצוע ניסוי בדיקה תוצאה מסקנה מבנה ארכיטקטורה תכנון יישום פיתוח הדמיה מיפוי ניווט מיקום זיהוי סיווג אומדן גישה מסגרת כלי תהליך שלב מרכיב רכיב מודול מנגנון ביצועים שיפור מחקר מערכת אלגוריתם ניתוח תוצאות שיטה מודל פתרון חישן נתון עיבוד מדידה הערכה ביצוע ניסוי בדיקה תוצאה מסקנה מבנה ארכיטקטורה תכנון יישום פיתוח הדמיה מיפוי ניווט מיקום זיהוי סיווג אומדן גישה מסגרת כלי תהליך שלב מרכיב רכיב מודול מנגנון ביצועים שיפור מחקר מערכת אלגוריתם ניתוח תוצאות שיטה מודל פתרון חישן נתון עיבוד מדידה הערכה ביצוע ניסוי בדיקה תוצאה מסקנה מבנה ארכיטקטורה תכנון יישום פיתוח הדמיה מיפוי ניווט מיקום זיהוי סיווג אומדן גישה מסגרת כלי תהליך

\en{See} \cite{Thrun2005}.
>>>>>>> d28438e746e968070f910361040b7f6540bdd144
